
\begin{abstract}
Routing in heterogeneous parallel queues requires a policy that responds to queue state, accounts for service heterogeneity, and remains stable. Classical dispatching rules address only part of this problem, while unguided policy-gradient methods offer little formal control. We introduce GibbsQ, an entropy-regularized routing framework that connects Lyapunov-stable analytical policies to teacher-initialized neural policies.

On the analytical side, we prove positive Harris recurrence for two routing laws---raw softmax and Unified Capacity-Aware Softmax (UAS)---under the strict load condition $\totcap > \arrate$. We also identify a certified subset within a broader calibrated extension, although the benchmark-default calibrated parameters do not satisfy the current sufficient condition. The drift constants are explicit, and for UAS they are independent of the concentration parameter.

On the empirical side, we introduce Reflected UAS, a three-parameter closed-form extension of UAS. Under the anchor benchmark, it reduces steady-state mean total queue length by 8.84\% relative to JSSQ and by 12.6\% relative to baseline UAS. These gains remain empirical rather than theorem-certified: a post-validation audit yields negative sufficient drift constants for the competitive reflected points we evaluate, although a benchmark-specific piecewise Lyapunov candidate suggests a plausible route toward future certification.

At the learned-policy layer, we isolate a teacher-choice effect. The default neural policy trained from the certified UAS teacher underperforms JSSQ and Reflected UAS, whereas a variant initialized from Reflected UAS and fine-tuned with \reinforce{} attains the lowest Protocol~B point estimate among all tested policies ($9.821$ versus $9.898$ for its closed-form teacher; paired $\Delta=-0.076$, 95\% CI $[-0.102,\,-0.051]$). Under the primary Protocol~A, the same variant matches Reflected UAS at approximately $10.04$. Because the Protocol~B result is the best of 5 random seeds selected on a held-out validation set, it should be interpreted as directional evidence rather than an unbiased ensemble expectation. All analysis and experiments assume standard M/M/1 dynamics; robustness to model misspecification is outside the present scope.
\end{abstract}
